\appendix
\section{Essential reproducibility details}
\label{app:essential}
\paragraph{Recognition data and optimization.}
Pets, CUB, Cars, and Dogs reserve development/calibration images within classes from the official training-side inventory. Flowers and Aircraft use official validation images for those roles. Exact and reviewed near-duplicate groups are separated from test-side images; additional cross-dataset exclusions precede mixed training. The duplicate audit covers the task-specific training and evaluation inventories. The six datasets contain 37, 200, 196, 120, 102, and 100 classes, respectively; the choice count is independent of class count.

FOCI/Finedefics test candidate sets are preserved, including the original repeated Dogs row, and shuffled consistently across methods. The unified protocol follows the released Dogs test inventory. Training and dynamic tests mine difficult incorrect labels with the recorded OpenCLIP ViT-L-14 image--text ranking; random distractors are distinct sampled incorrect classes. The correct label is included. The standard author-defined four-choice test and the newly constructed difficult four-choice test are separate conditions.

Unified training uses learning rate 0.0001, accumulation over eight questions, three epochs, cosine decay, warmup fraction 0.03, AdamW weight decay 0.01, global clipping at norm one, and language-attention LoRA rank/scaling 16/32. The seed is 20261021. Both methods use candidate-conditioned cross-entropy, the same exposure and selection rule, and development-based checkpoint selection. Variable-count training changes only the registered candidate-count sampling and associated candidate construction; candidate-token computation follows the sampled candidate counts. Images use processor area bounds 65536 and 262144 pixels; Aircraft crops the bottom twenty pixels. Other recognition images are uncropped.

\paragraph{Relation construction and endpoint.}
Images contain colored objects with distinct sizes. Paired three-choice questions share two middle objects but replace the third candidate with a smaller or larger extreme, reversing the median answer. The largest-object task serves as an independent-ranking control.  Whole-scene bootstrap intervals condition on those trained seeds.

\paragraph{VQA comparison scope.}
The main text presents the registered native-LM secondary comparison. The original primary contrast used a canonically sorted LM input: Set-minus-canonical-LM accuracy was $-0.5694$ percentage points, with conditional interval $[-1.2861,0.1289]$. The original primary and secondary analysis roles are retained. For native LM, raw Set-minus-LM NLL and Brier differences are $+0.1933$ and $+0.0243$; their exploratory crossed intervals are $[0.0779,0.3260]$ and $[0.0078,0.0407]$. Scaled differences are $+0.0116$ for NLL and $+0.0087$ for Brier, with crossed intervals $[-0.0054,0.0284]$ and $[-0.0025,0.0204]$. The registered analysis estimates accuracy and probability-loss differences with intervals.

\paragraph{Calibration and statistical conventions.}
Recognition minimizes reserved-set pooled candidate NLL over $\log T\in[-5,5]$ using bounded scalar optimization and explicit boundary checks. VQA uses a fixed grid of 2001 log-temperatures in the same range, preferring the value closest to zero on objective ties. Temperatures are fixed before test scoring. Sum-Brier sums squared probability errors across candidate or literal-answer events. Recognition ECE uses fifteen equal-count bins; VQA uses fifteen equal-width bins. Exact score ties are resolved by selecting the first slot. Paired bootstrap samples keep repeated images and related questions together. The recognition and relation intervals condition on the trained seeds; the crossed VQA analysis is exploratory with few training realizations. All intervals are conditional on the fitted temperatures and the evaluation distribution.

\paragraph{Reproduction scope.}
The recorded implementation uses PyTorch 2.8.0, Transformers 4.57.3, PEFT 0.17.1, and NVIDIA A800 devices. Data manifests, model checkpoints, raw predictions, and analysis scripts provide the evidence trail. The project archive preserves earlier specialist adapters, intermediate training endpoints, optimization probes, and engineering validation records. Public external reproduction remains a planned release milestone.

\subsection{Equivalent answer events}
For the registered VQA protocol, blank and duplicate candidate strings are preserved. Prediction selects the maximum-logit slot, with first-slot tie breaking, before comparison by literal answer text. For probability scoring, the mass of identical strings is summed. Given a gold answer event $G$, the candidate-conditioned loss is
\begin{equation}
 \mathcal L=-\log\sum_{i\in G}p_i
 =\operatorname{LSE}_{j\in\mathcal V}(s_j)-\operatorname{LSE}_{i\in G}(s_i).
\end{equation}
Labels supply supervision to the loss and evaluator; forward prompts contain the image, question, and candidate text. This convention preserves the official answer-text task while keeping slot and answer-event probabilities distinct.

\subsection{Supplementary evaluation details}
\paragraph{Candidate-count study.}
The variable-count training ablation uses a separate unified CE recipe: three epochs with matched image exposure, alternating difficult and random distractors, and development selection by equal-dataset accuracy with earlier-epoch tie-breaking. Fixed-count training uses four candidates; variable-count training samples uniformly from $K\in\{2,4,10\}$. The reported count-transfer tests use difficult distractors at two, four, ten, and twenty candidates alongside the standard records. The four checkpoints share the training seed and test manifests. This study matches image exposure and updates; candidate-token computation varies with candidate count. Twenty-candidate testing measures transfer beyond the training count range.

\paragraph{Order and relation tests.}
The order test evaluates all 24 permutations of four candidates on 755 class-stratified questions, aligning candidate identities before comparison. Both this test and the difficult-candidate count-transfer evaluation use the complete-method checkpoints from Table~\ref{tab:unified-standard}, with unchanged weights and inference precision. Relation tests use synthetic scenes with paired candidate sets that reverse the required preference between shared objects. Independent, Set, and complete-candidate LM are evaluated with three training seeds at the same fixed final training endpoint on fresh scenes. The paired metric requires both questions correct; a largest-object control tests ordinary independent ranking.

\paragraph{Probability quality and uncertainty.}
Each model receives its own scalar temperature fitted only on reserved calibration data~\citep{guo_2017_calibration}. We report accuracy, NLL, sum-Brier, and ECE, distinguishing raw and scaled probabilities. Recognition ECE uses fifteen equal-count bins; the VQA protocol uses equal-width bins over literal answer events. Recognition intervals use 10,000 paired image-cluster bootstrap replicates within datasets and equal dataset weighting; repeated image rows stay together. They condition on the selected weights and fitted temperatures, excluding training-seed and calibration-fit uncertainty. Relation intervals resample whole scenes; the VQA interval additionally resamples the few available seeds and is exploratory. Intervals are marginal, without multiplicity correction. Further protocol details appear in Appendix~\ref{app:essential}.
